% !PRINTED 574
\weberpage{574}
If we omit from $\Phi(t)$ and $\Phi_0(t)$ the terms whose coefficients are divisible by $p$, then the degree of $\Phi(t)$ with respect to $t$ is greater by $f$ than the degree of $\Phi_0(t)$.

$P(t)$ is an integral function in $R$ of the variables $t, t_1, t_2, \ldots, t_n$, which by the substitution $t=\tau$ goes over into a functional divisible by $\mathfrak p$, which of course is no longer contained in $R$.

The significance of this form $P(t)$ now emerges still more clearly when we prove the theorem:

4. \gesperrt{The prime factor $\mathfrak p$ is the greatest common divisor of $p$ and $P(\tau)$.}

For this we have only to show that, if $p$ is divisible by $\mathfrak p \mathfrak p_1$, where $\mathfrak p, \mathfrak p_1$ are two equal or different prime factors, then $P(\tau)$ is divisible by $\mathfrak p$, but not by $\mathfrak p \mathfrak p_1$.

By \S.\,143, 3. there always exists a \gesperrt{number} $\xi$ in $\mathfrak{o}$ which is divisible by $\mathfrak p$, but not by $\mathfrak p \mathfrak p_1$. This number will have the form
\weq{16}{\xi = a_1 \omega_1 + a_2 \omega_2 + \cdots + a_n \omega_n,}
wherein the $a_1, a_2, \ldots, a_n$ are integral rational numbers, and it therefore arises from the form $\tau$ by the substitution
\weq{17}{(t_1, t_2, \ldots, t_n) = (a_1, a_2, \ldots, a_n).}
If we make the same substitution in the forms $\tau_r$ defined in (9), then $\tau_r$ goes over into a number $\xi_r$ which, by \gesperrt{Fermat}'s theorem, satisfies the congruence
\[ \xi^{p^r} \equiv \xi_r \pmod{\mathfrak p} \]
and is thus certainly also divisible by $\mathfrak p$. If, therefore, in the form
\[ \Pi(t) = (t-\tau)(t-\tau_1) \cdots (t - \tau_{f-1}) \]
we put $\tau_r + \xi_r$ in place of $\tau_r$, then this form remains congruent to itself with respect to the modulus $\mathfrak p$. But this substitution comes to the same thing as substituting $t_1+a_1, t_2+a_2, \ldots, t_n+a_n$ in place of $t_1, t_2, \ldots, t_n$, and we accordingly obtain
\[ \Pi(t, t_1+a_1, \ldots, t_n+a_n) \equiv \Pi(t, t_1, \ldots, t_n) \pmod{\mathfrak p}, \]
% [sic: Weber prints ``nnd'' for ``und'' (and).]
and if we replace $\Pi$ by the congruent form $P$
\[ P(t, t_1+a_1, \ldots, t_n+a_n) \equiv P(t, t_1, \ldots, t_n) \pmod{\mathfrak p}. \]
